%% Section 3: background on graph models. Seed: the figure of elements and
%% operators that the later sections use. Budget: 2.0 pages
%% (notes/paper-outline.md). Definitions follow the Agda:
%% agda/Framework/GraphModel/Domain.agda (elements, consistency),
%% agda/Framework/GraphModel/Ops.agda (operators),
%% agda/proof/Framework/GraphModel/DomainUtilBase.agda (projections, depth).
%% Exposition after Siek 2017 (Revisiting Elementary Denotational
%% Semantics, cited in the third person) and the Denotational chapter of
%% PLFA: functions as finite tables, self-application by inclusion,
%% divergence as the empty set. Line counts: notes/paper-outline.md, section 3.

\section{Graph Models}
\label{sec:graph-models}

\begin{figure}[t]
\small\figdefs
\defhead{Graph elements}{u, v, w \in \Elem}
\vspace{-10pt}
\[
\begin{array}{@{}r@{\;}c@{\;}l@{}}
  u, v, w & ::= & c \mid V \mapsto w \mid \nu \mid \pairv{u}{v}
    \mid \mathsf{tup}_i(v) \mid \mathsf{left}(v) \mid \mathsf{right}(v)
    \mid \langle\rangle
\end{array}
\]
\defhead{Finite lists of elements}{U, V, W \in \mathsf{List}(\Elem)}
\defhead{Denotations (sets of elements)}{D, E \subseteq \Elem}
\defhead{Indices}{i, j \in \mathbb{N}}
\defhead{Consistency (Appendix~\ref{app:consistency})}{u \sim v \qquad \mathsf{con}(V)}
\vspace{-10pt}
\[
  \mathsf{con}(V) \iff u \sim v \text{ for all } u, v \in V
\]
\defhead{Operators on denotations}{\mathit{op}(D_1, \ldots, D_n) \subseteq \Elem}
\[
\begin{array}{@{}r@{\;}c@{\;}l@{}}
  \Lambda(F) & = & \{\, V \mapsto w \mid V \neq [\,],\; \mathsf{con}(V),\; w \in F(V) \,\} \cup \{\nu\}
    \quad (F \text{ a function on sets, e.g. } \mlam{X} E)\\
  D \app E & = & \{\, w \mid V \mapsto w \in D,\; V \neq [\,],\; V \subseteq E \,\}\\
  \mathsf{pair}(D, E) & = & \{\, \pairv{u}{v} \mid u \in D,\; v \in E \,\}\\
  \car(D) & = & \{\, u \mid \pairv{u}{v} \in D \,\}\\
  \cdr(D) & = & \{\, v \mid \pairv{u}{v} \in D \,\}\\
  \mathcal{T}(D_1, \ldots, D_n) & = &
    \{\, \mathsf{tup}_i(v) \mid v \in D_i \,\}
    \quad \text{if } D_1, \ldots, D_n \neq \emptyset,
    \text{ and } \emptyset \text{ otherwise}\\
  \mathsf{inl}(D) & = & \{\, \mathsf{left}(v) \mid v \in D \,\}
    \qquad
  \mathsf{inr}(D) \;=\; \{\, \mathsf{right}(v) \mid v \in D \,\}\\
  \pi_t(D) & = & \{\, v \mid t(v) \in D \,\}
    \quad \text{for } t \in \{\mathsf{tup}_i, \mathsf{left}, \mathsf{right}\}\\
  \mathcal{C}(D, F, G) & = &
    \{\, w \mid V \neq [\,],\; \mathsf{con}(V),\; \mathsf{left}(V) \subseteq D,\; w \in F(V) \,\}\\
  & \cup & \{\, w \mid V \neq [\,],\; \mathsf{con}(V),\; \mathsf{right}(V) \subseteq D,\; w \in G(V) \,\}\\
  \mathsf{left}(V) & = & \{\, \mathsf{left}(v) \mid v \in V \,\}
    \qquad
  \mathsf{right}(V) \;=\; \{\, \mathsf{right}(v) \mid v \in V \,\}\\
  \mathsf{prim}_p(D, E) & = & \{\, \delta(p, n, m) \mid n \in D,\; m \in E \text{ integers} \,\}\\
  \delta(\kw{sub}, n, m) & = & n - m
    \qquad
  \delta(\kw{eq}, n, m) \;=\;
    \mathsf{left}(\langle\rangle) \text{ if } n = m,
    \text{ and } \mathsf{right}(\langle\rangle) \text{ otherwise}
\end{array}
\]

\defhead{Semantics of \fpc}{\den{M}\rho = D}
\[
\begin{array}{@{}r@{\;}c@{\;}l@{\qquad\qquad}r@{\;}c@{\;}l@{}}
  \den{x}\rho & = & \rho(x) &
  \den{c}\rho & = & \{c\}\\
  \den{\lambda x.\, N}\rho & = & \Lambda(\mlam{X} \den{N}\rho[x{:=}X]) &
  \den{L\;M}\rho & = & \den{L}\rho \app \den{M}\rho\\
  \den{\kw{tuple}\;(M_1, \ldots, M_n)}\rho & = & \mathcal{T}(\den{M_1}\rho, \ldots, \den{M_n}\rho) &
  \den{\kw{get}_i\,M}\rho & = & \pi_{\mathsf{tup}_i}(\den{M}\rho)\\
  \den{\kw{inl}\;M}\rho & = & \mathsf{inl}(\den{M}\rho) &
  \den{\kw{inr}\;M}\rho & = & \mathsf{inr}(\den{M}\rho)\\
  \den{\kw{prim}_p\;L\;M}\rho & = & \mathsf{prim}_p(\den{L}\rho, \den{M}\rho) &
  \den{\kw{fold}\;M}\rho & = & \den{M}\rho\\
  \den{\kw{let}\;x = M\;\kw{in}\;N}\rho & = & \Lambda(\mlam{X} \den{N}\rho[x{:=}X]) \app \den{M}\rho &
  \den{\kw{unfold}\;M}\rho & = & \den{M}\rho
\end{array}
\]
\[
  \den{\kw{case}\;L\;(x.\,M)\;(y.\,N)}\rho =
    \mathcal{C}(\den{L}\rho,\; \mlam{X} \den{M}\rho[x{:=}X],\; \mlam{Y} \den{N}\rho[y{:=}Y])
\]
\caption{The graph model (elements, denotations, and operators) and
  the semantics of \fpc.}
\label{fig:graph-ops}
\end{figure}

A denotational semantics maps each program to a mathematical object
that describes its behavior. In a graph model, that object is a set of
finite pieces of evidence about the behavior, which we call
\emph{elements}. Our model is a variant of the model of
\citet{Engeler:1981aa}, which we extend with constants, pairs, tuples,
and disjoint sums. Figure~\ref{fig:graph-ops} defines elements,
several semantic operators on sets of elements, and the semantics of
\fpc. We first explain how the model represents functions, then
introduce the elements, and then the operators. We review how the
model handles recursion and divergence, and why we require sets of
elements to be consistent. We end this section with the properties of
the semantics that our proofs rely on, including a $\beta$-law.

\paragraph{Functions as tables}

One can describe a function by tabulating it. Each row of the table
pairs an input with the corresponding output, and applying the
function means looking up a row. Two problems arise with this simple
approach. First, a function on an infinite domain (such as the
integers) needs an infinite table. Second, a function that is applied
to itself (such as the Y combinator) would need a table that contains
a copy of itself. Graph models solve both problems. For the first
problem, a terminating program applies each function to only finitely
many arguments, so finite tables suffice
\citep{Scott:1976lq,Plotkin:1993ab}. However, to handle all possible
contexts, the meaning of a function is the set of all its finite
tables, or equivalently the set of all their rows. For the second
problem, a row applies to an argument when the row's input is
\emph{included} in the argument, rather than equal to it. A table then
needs only a smaller copy of itself as an input entry, and elements,
defined as an inductive data type, can provide that.

\paragraph{Graph elements}

An element is a finite and partial description of a term's behavior
(Figure~\ref{fig:graph-ops}). A row of a table is represented by an
element $V \mapsto w$, called an \emph{entry}, whose input $V$ is a
list of elements\footnote{We sometimes use a list of elements where a
set of elements is expected, and assume the natural conversion takes
place.}  and whose output $w$ is an element.\footnote{The input of an
entry could instead be a single element, as in filter models
\citep{Barendregt:1983aa}, but a list makes the metatheory cleaner
because the argument of a function is often a function, and $V$ can
contain one entry for each call to it.} The element $\nu$ records that
a term evaluates to a function, without describing any of the
function's rows. The other elements describe the data of our
languages: a constant $c$, a pair $\pairv{u}{v}$, a tuple element
$\mathsf{tup}_i(v)$ which says that component $i$ of a tuple is
described by $v$, the empty tuple $\langle\rangle$, and the injections
$\mathsf{left}(v)$ and $\mathsf{right}(v)$ of a sum. A denotation is a
set of elements.

\paragraph{Operators}

Each case of the semantics of \fpc\ (Figure~\ref{fig:graph-ops})
applies an operator to the denotations of the sub terms. In
the semantics, an environment $\rho$ maps each variable to a set of
elements, $\rho[x{:=}D]$ maps $x$ to $D$ and every other variable as
$\rho$ does, and $\boldsymbol{\lambda}$ denotes a mathematical
function.

The meaning of a function, given by the operator $\Lambda$, is the set
of all entries $V \mapsto w$ such that the output $w$ is in the meaning of
the function's body in an environment that maps its parameter to the
input $V$, together with $\nu$. The input $V$ must be nonempty and
consistent, written $\mathsf{con}(V)$, which we explain at the end of
this section. Because of $\nu$, a $\lambda$-abstraction has a nonempty
meaning, as befits a value in a call-by-value language, even if it is
never applied. The operator $D \app E$ gives the meaning of
application. It collects the outputs of the entries of $D$ whose input
is contained in the argument $E$. For example, $\den{\lambda x.\, x +
  1}$ contains the entry $[n] \mapsto n{+}1$ for every integer $n$, and
$\den{(\lambda x.\, x + 1)\; 3} = \{4\}$.
%
For a higher-order example, $\den{\lambda f.\, \lambda x.\, f\;(f\;x)}$
contains the entry
\[
[\, [0] \mapsto 1,\; [1] \mapsto 2 \,] \mapsto
([0] \mapsto 2)
\]
whose input has one entry for each call to $f$. That input is included
in $\den{\lambda x.\, x + 1}$, so applying the higher-order function
to $\lambda x.\, x + 1$ and then to $0$ gives $\{2\}$.

The other operators build and take apart data. The operator
$\mathsf{pair}$ builds pairs, and $\car$ and $\cdr$ take their
components. The operator $\mathcal{T}$ builds tuples, and
$\pi_{\mathsf{tup}_i}$ takes component $i$. The operators
$\mathsf{inl}$ and $\mathsf{inr}$ build injections, and $\mathcal{C}$
gives the meaning of a $\kw{case}$. Like $\Lambda$, it binds the
variable of a branch to a nonempty, consistent list, here of the
payloads of the injections in the scrutinee. The operator
$\mathsf{prim}_p$ applies a primitive to the integers in its operands.
The empty tuple $\langle\rangle$ also fills the positions of an
environment that hold no variable, because we represent environments
as total functions whose positions are nonempty. The environment
$\rho_0$ of a closed program maps every variable to
$\{\langle\rangle\}$, and $\den{M} = \den{M}\rho_0$.

\paragraph{Recursion and divergence}

Recursion in \fpc\ comes from self-application, which the recursive
types of \fpc\ make typable. The constructs $\kw{fold}$ and
$\kw{unfold}$ mean the identity, so recursive types cost nothing in
the semantics. Because entries apply to arguments that include their
input, the self-applied function in a fixed-point combinator such as
$Z$ has a meaning whose tables contain smaller copies of themselves,
like a set of nested dolls \citep{Siek:2017ab}.

A program that diverges produces no evidence, so its meaning is the
empty set. For example, $\Omega = \Delta\;(\kw{fold}\;\Delta)$ with
$\Delta = \lambda x.\, (\kw{unfold}\;x)\; x$ denotes $\emptyset$, because an
output of $\Omega$ would need an entry of $\den{\Delta}$ whose input
contains a strictly smaller entry of $\den{\Delta}$ with the same
output, and so on forever, which is impossible for finite elements.


\paragraph{Consistency}

We write $u \sim v$ when the elements $u$ and $v$ are consistent, which means
that they could describe the same value. Appendix~\ref{app:consistency} gives
the definition. The operator $\Lambda$ builds only entries whose inputs are
consistent, which rules out entries such as $[1, 2] \mapsto 3$ in the meaning of
$\lambda x.\, x + 1$, because the input $[1, 2]$ of that entry
describes no single argument.
%
Our model requires consistency because the proof for the \pass{split closure} pass needs it (Section~\ref{sec:approx-lr}). A call in PairClos takes
the code from one pair and the environment from another, so a
closure's meaning may combine the code of one pair with the
environment of a different pair. Consistency ensures that the
environments of these pairs are consistent with each other.


\paragraph{The properties of the semantics}

%% Lemma: agda/proof/Framework/GraphModel/Ops.agda, Λ-β (and ΛΛ-β for
%% the curried form used in the delay proof).
%% The finitary premise holds for term meanings by `ok` in
%% agda/FPC/proof/Semantics/GraphModel/{FPC,EnvClos,PairClos,GlobalFun}.agda (fin₁, fin₂).
The meaning of a term is monotone in its environment
(Lemma~\ref{lem:monotone}), and it is continuous in the sense that
each element of the meaning depends on only finitely many elements of
the environment. These properties, together with consistency, give the
graph model the following $\beta$ laws. Its premises come from the
definitions of $\Lambda$ and application, which use only entries whose
input is nonempty and consistent.
%
A function $F$ on sets is called \emph{finitary} at $D$ if every $w
\in F(D)$ is in $F(V)$ for some nonempty $V \subseteq D$, and
similarly for a function of two sets.


\begin{lemma}[$\beta$ laws]\label{lem:beta}\ 
\begin{enumerate}
\item If $F$ is monotone and finitary at $D$ and $D$ is consistent,
  then \\
  $\Lambda(F) \app D = F(D)$.
\item If $F$ is finitary at $T$ and $E$, and both $T$ and $E$ are consistent,
  then \\
  $F(T, E) \subseteq (\Lambda(\mlam{X} \Lambda(\mlam{Y} F(X, Y))) \app T) \app E$.
\end{enumerate}
\end{lemma}
